% !TEX root = ../../main.tex
\subsection*{Biological sliding-window analyses}

The protein 2D and rRNA 3D examples use the alignments and topologies of Hug and colleagues \cite{hug2016}. Branch lengths are re-estimated in IQ-TREE under LG+G4 for the 2,596-site, 3,083-taxon protein alignment and GTR+G4 for the 1,947-site, 1,871-taxon rRNA alignment. Each site is assigned to the rate category with largest posterior probability. Whole-category branch tests use the taxon-pair Bonferroni threshold because a category contains roughly one quarter of the alignment.

Sliding windows contain 36 sites, the minimum scale used in the original analysis. A global variance estimate is formed from the category-specific estimates weighted by their full-alignment site fractions; each window uses its mean coherence and this global variance. The original analysis does not apply Bonferroni correction to windows because one window contains less than 2\% of the alignment used to infer the tree. Saturated windows satisfy \(z\leq z_\alpha\). Hypervariable rRNA regions V1--V9 and conserved regions C1--C9 are mapped from bacterial 16S annotations after accounting for alignment gaps \cite{baker2004,yang2016}.

The curated rRNA rerun compares published SatuTe with native recomputations of the published and eigenvalue-weighted SatuTe statistics for the Eukaryota and yeast branches. The independent rerun reconstructs site scores from the IQ-TREE treefile, model report, site-probability output and original alignments for both Tree-of-Life datasets. Cross-implementation count differences are reported descriptively and are not treated as paired effects of eigenvalue weighting.

For the topology-rearrangement analysis, the Eukaryota subtree is pruned from the protein 2D tree and regrafted onto the Bacteria--Archaea branch to create a 3D alternative; the reverse operation creates a 2D alternative from the rRNA 3D tree. Branch lengths are re-estimated under the corresponding LG+G4 or GTR+G4 model. Protein-specific scores are calculated for each of 16 concatenated proteins, and rRNA-specific scores for nine annotated regions. Differences are defined consistently between 2D and 3D placements and compared descriptively, with a two-sample Kolmogorov--Smirnov test for the two empirical difference distributions.

For the HCV example, complete accession-versioned references were selected for subtypes 1a, 1b, 2a, 2b, 3a, 3b, 4a, 5a, 6a and 7a. Polyprotein translations were aligned with MAFFT 7.526 and back-translated; H77 mature-peptide annotations defined Core--NS5B. A full-coding topology was inferred under GTR+F+G4, and the subtype pairs for genotypes 1--3 defined target branches. Codon positions 1+2 and position 3 were concatenated separately. Fixed-topology 36-site windows advanced by 12 sites, with model parameters and branch lengths re-estimated per window. Positions 1+2 and position 3 therefore had equal statistical sample size but spanned 18 and 36 aligned codons.

For every HCV window we recorded the standardized published and eigenvalue-weighted SatuTe statistics and plotted them at the corresponding H77 polyprotein midpoint. The horizontal decision line is \(z_{0.01}=2.326\); a window is counted as saturated when its \(z\)-score falls below this line. The 507 positions-1+2 windows and 252 position-3 windows were required for each of three target branches and both statistics, yielding a prespecified 4,554-row grid. We did not correct across overlapping windows because these curves are descriptive localization profiles. The ten-reference design tests partition localization, not population diversity, recombination or genotype assignment.
